\documentclass[runningheads]{llncs}
\usepackage[T1]{fontenc}
\usepackage{lmodern}
\usepackage{amsmath,amssymb,booktabs,array,microtype}
\usepackage[hidelinks]{hyperref}
\usepackage{xurl}
\begin{document}
\title{When Set Scores Fail to Select Evidence:\\A Controlled Study of Budgeted Multi-hop RAG}
\titlerunning{When Set Scores Fail to Select Evidence}
\author{Anonymous authors}
\authorrunning{Anonymous authors}
\institute{Anonymous affiliation}
\maketitle
\begin{abstract}
Selecting complementary evidence under a limited context budget requires a
score that remains useful on the sets found by search. We examine this
requirement with first-, second-, and fourth-order factorization models and
a permutation-invariant set model, sharing dense candidates and
support-annotation supervision. Strong performance on constructed set pairs
does not translate into stronger evidence selection: the static fourth-order
model loses complete annotated support relative to a size-matched dense
baseline. A bounded correction trains the deployed score on actual search
frontiers and final sets, with an equal-update static replay control.
For the primary seed, fourth-order checkpoint-panel completeness increases
from 0.4531 to 0.5156, but answer F1 on a separate 128-question development
panel falls from 0.4181 to 0.3766, below maximal marginal relevance's 0.4303.
In both training seeds, every aligned model remains below this baseline
in F1 point estimate.
Exact granular-ball and k-means indexes preserve flat-search
selections but do not accelerate the measured 128-candidate search.
The study distinguishes proxy fit, selection quality, answer quality, and
runtime, identifying where additional model structure does and does not
provide practical value. All learned-selector comparisons are development
evidence, rather than a new independent confirmation.
\keywords{Evidence selection \and Multi-hop QA \and Search alignment}
\end{abstract}

\section{Introduction}
Multi-hop question answering often depends on a small combination of
passages rather than the single passage most similar to the question.
Under a limited reader budget, evidence selection must balance relevance,
redundancy, and coverage. Set-wise selection makes this collective objective
explicit~\cite{lee2025setr}, while beam retrieval preserves multiple partial
evidence hypotheses~\cite{zhang2024beam}. Neither modeling interactions nor
searching more combinations ensures that the resulting context will help a
reader.

We investigate a particular failure mode: a learned evidence-set score can
perform well on constructed training comparisons while assigning high
scores to poorer sets encountered during deployment. This motivates a
controlled study with frozen embeddings, shared candidates, a fixed
reader, and four comparable set models. The first-, second-, and
fourth-order models test the value of explicit interactions; a DeepSets
control tests whether any benefit requires their particular factorization.
Earlier HyperGranular-RAG sentence-completion experiments motivate this
distinct paragraph-selector study.

Our contribution is an empirical separation of three gaps: constructed-set
prediction versus actual selection, support coverage versus answer quality,
and fewer score computations versus lower runtime. A bounded search-aligned
training intervention probes the first gap without changing model families
or first-stage retrieval. It follows established search-aware learning
ideas~\cite{wiseman2016bso}; we claim neither a new factorization algebra nor a
new tree bound. This framing makes the simple controls and negative cost
results central to the research question.

\section{Budgeted Evidence Selection}
\paragraph{Shared task.}
For question $q$, BGE-large-en-v1.5~\cite{xiao2023bge} retrieves a common candidate set $V_q$ of at most 128
whole paragraphs from a corpus assembled from exposed development data. A context contains
$S\subseteq V_q$, $|S|\leq6$, rendered in original BGE order. The actual
tokenized reader input, including the question and instructions, must not
exceed 1,024 tokens. This token constraint is distinct from the six-block
constraint. No support paragraph is inserted into evaluation candidates.

Let $G_q$ denote the known support annotations and $C(S)$ the annotations
covered by the selected original passages. Supervision is
\begin{equation}
 c(S)=|C(S)\cap G_q|/|G_q|,\qquad f(S)=\mathbf1[G_q\subseteq C(S)].
\end{equation}
These labels measure annotated support, not factual truth or exact answer
utility. Unannotated passages may still be useful. Training can use known
support references, whereas deployment selection receives only the question,
passages, embeddings, and token lengths.

\paragraph{Set scores and search.}
Query--passage features combine frozen normalized embeddings, their product
and absolute difference, cosine similarity, and log length. A shared
128-dimensional network produces unary factors $a_i$ and 32-dimensional
interaction factors $v_i$. For order $H\in\{1,2,4\}$, the deployed score is
\begin{equation}\label{eq:score}
g_H(S)=b_q(|S|)+\frac1{6}\sum_{i\in S}a_i+
\sum_{h=2}^{H}\frac{w_{q,h}^{\top}e_h(\{v_i:i\in S\})}{32\binom6h},
\end{equation}
where $e_h$ is the coordinatewise elementary symmetric polynomial.
Its recurrence is the standard distinct-member ANOVA computation from
higher-order factorization machines~\cite{blondel2016hofm}.
The DeepSets control replaces the interaction sum with an MLP of the summed
passage features and query representation~\cite{zaheer2017deepsets}.
Each model has a full-support logit and an auxiliary coverage head.

Width-four beam search evaluates feasible extensions through depth six,
retains visited sets, and includes a feasible Dense-K6 set in the final
comparison. It returns the highest-scoring visited set, with token length
then stable IDs resolving final ties. Beam search is not a global optimum.
Dense-budget and MMR-budget may retain more than six short blocks under the
same token limit; Dense-K6 and MMR-K6 isolate cardinality in support analyses.
MMR uses the fixed score $0.7\,s(q,i)-0.3\max_{j\in S}s(i,j)$.

\paragraph{Search-aligned continuation.}
Static supervision contains support sets, missing-support variants,
same-size replacements, redundant additions, and random sets. We continue
each model on a shared pool mined from actual FIT search: model final sets,
partial frontiers, Dense-K6/MMR-K6, and limited replacements. Pools are
merged across all four architectures, deduplicated, and capped at 64 sets
per question. Reference support can augment static FIT examples but never
the mined candidate universe.

For $u(S)=(f(S)+c(S))/2$ and $u(S^+)>u(S^-)$, the added preference loss is
\begin{equation}
\ell(S^+,S^-)=\operatorname{softplus}\!\left(
0.2[u(S^+)-u(S^-)]-[g(S^+)-g(S^-)]\right).
\end{equation}
Final pairs prioritize equal cardinality and similar length; frontier pairs
prioritize the same parent and depth. Thus coverage differences between two
incomplete sets supervise the score used by search directly. The objective
is $L_{\mathrm{base}}+0.5L_{\mathrm{final}}+0.5L_{\mathrm{frontier}}$,
with query-normalized losses and equal static/mined base sampling.
H4 static replay receives the same optimizer settings and number of updates,
without mined examples or the added preferences. Search-aware training and
overestimation control have established precedents
\cite{wiseman2016bso,trabucco2021com}; our intervention tests their relevance
to this annotation-supervised selection problem, without borrowing their
performance or theoretical guarantees.

\section{Experimental Setting}
We reuse previously exposed HotpotQA and MuSiQue materials
\cite{yang2018hotpotqa,trivedi2022musique}. Paragraph reconstruction yields
9,791 and 31,130 blocks. The 4,000 questions are divided into 3,166
FIT, 738 TUNE, and 96 historical-development questions. Query groups remain
separate, but FIT and TUNE share 57 HotpotQA and 4,753 MuSiQue documents.
This is development-corpus retrieval, not full-Wikipedia or unseen-document
evaluation. The 738-query selection analysis contains 196 HotpotQA and 542
MuSiQue questions. Top128 reaches all known support for 190 and 367,
respectively. Nine MuSiQue questions have known support that cannot all fit
the six-block/1,024-token constraint. All remain in the denominator.

The continuation mines 256 FIT queries per dataset and selects checkpoints
using 64 non-QA TUNE queries per dataset. Hash-based sampling preserves
existing question-type proportions and query groups. Up to two eight-epoch
rounds use AdamW (learning rate $3\times10^{-4}$, weight decay $10^{-4}$,
batch eight queries, gradient norm capped at one). At epochs 2, 4, and 8,
checkpoint selection uses dataset-equal actual support completeness, then
coverage, original development loss, and earlier epoch. QA never selects
the checkpoint. The replay checkpoint follows the selected H4 update count.
Both seeds (1729/2026) complete two rounds; the second seed was triggered
by selection-panel improvement before new QA scores were observed.

The fixed QA panel contains 64 questions per dataset: 48 bridge and 16
comparison HotpotQA questions, and 32 two-hop and 32 longer-hop MuSiQue
questions. The original FP16 Qwen2.5-3B-Instruct reader~\cite{qwen2024release} uses greedy decoding,
at most 128 output tokens, and a shared constrained single-answer-field
protocol. Dataset-specific canonical EM/F1 score that field without
answer-dependent cleaning. Identical full token/model/decoder identities
reuse outputs; this is not a new generation. Failures remain in every mean.
All panel questions are exposed development data. We report paired
descriptive comparisons without new confirmation tests or significance
claims. Three original training seeds are not three independent datasets.

\section{Results}
\paragraph{Good proxy fit can coexist with support loss.}
Table~\ref{tab:proxy} contrasts constructed-pair accuracy with actual
selection. H4 correctly orders most sampled complete/incomplete pairs, yet
selects less complete support than Dense-K6 on both datasets. Relative to
Dense-K6, it loses 27 and 42 complete-support queries while recovering only
5 and 24. Its mean sigmoid outputs on selected sets are 0.992 and 0.967,
against actual completeness 0.709 and 0.223. These outputs are not calibrated
probabilities. The discrepancy motivates training on states produced by
search rather than improving only an auxiliary prediction metric.

\begin{table}[t]
\caption{Original static models (seed 1729): constructed-pair accuracy and actual complete
support on all 738 TUNE questions. H/M denote HotpotQA/MuSiQue.}
\label{tab:proxy}\centering\small
\begin{tabular}{lrrrr}\toprule
Method & Pair H & Pair M & Full H & Full M\\\midrule
Dense-K6 & -- & -- & .8214 & .2565\\
H1 & .9475 & .8978 & .6531 & .2454\\
H2 & .9638 & .9422 & .6429 & .2325\\
DeepSets & .9620 & .9482 & .6990 & .2306\\
H4 & .9645 & .9464 & .7092 & .2232\\\bottomrule
\end{tabular}
\end{table}

\begin{table}[t]\centering\small
\caption{Checkpoint-panel support completeness before and after search alignment (64 questions per dataset). Selection uses this panel, not QA.}
\label{tab:aligned}
\begin{tabular}{llrrrr}\toprule
Seed & Model & Static H & Aligned H & Static M & Aligned M\\\midrule
1729 & H1 & 0.5938 & 0.7500 & 0.2656 & 0.2656\\
1729 & H2 & 0.6406 & 0.7500 & 0.2344 & 0.2500\\
1729 & DeepSets & 0.6719 & 0.7500 & 0.2188 & 0.2656\\
1729 & H4 & 0.6719 & 0.7500 & 0.2344 & 0.2812\\
2026 & H1 & 0.6094 & 0.7656 & 0.2344 & 0.2656\\
2026 & H2 & 0.6406 & 0.7188 & 0.2500 & 0.2500\\
2026 & DeepSets & 0.6250 & 0.7812 & 0.2344 & 0.2656\\
2026 & H4 & 0.6719 & 0.7188 & 0.2344 & 0.2500\\
\bottomrule\end{tabular}\end{table}

Table~\ref{tab:aligned} shows actual checkpoint-panel selections. For seed 1729, H1 changes equal-dataset completeness by $+0.0781$, H2 changes equal-dataset completeness by $+0.0625$, DeepSets changes equal-dataset completeness by $+0.0625$, H4 changes equal-dataset completeness by $+0.0625$. These are selected development effects, with the full checkpoint trajectory retained. They do not by themselves establish answer improvement or a unique advantage of fourth-order interactions. Every reported aligned checkpoint comes from round one; the completed second rounds did not improve the best selection criterion. Continued loss reduction was not continued selection progress.


\paragraph{Answer quality remains a separate endpoint.}
Table~\ref{tab:qa} retains Dense, MMR, the original H4 checkpoint, replay,
and every aligned architecture on the same 128 questions. In the earlier
three-seed experiment, H4's mean dataset-equal F1 was 0.4171 (seed standard
deviation 0.0139), compared with Dense's 0.4055 and MMR's 0.4303.
The H4 range is 0.4027--0.4304; seed variation is neither a query confidence
interval nor an ensemble result.
\begin{table}[t]\centering\small
\caption{Canonical answer F1 on the fixed 128-query development panel (64 per dataset), at 1,024 input tokens. Original rows reuse their exact prior generations.}
\label{tab:qa}
\begin{tabular}{llrrr}\toprule
Seed & Method & HotpotQA & MuSiQue & Mean\\\midrule
1729 & Dense & 0.6324 & 0.1786 & 0.4055\\
1729 & MMR & 0.6615 & 0.1991 & 0.4303\\
1729 & Static H4 & 0.6341 & 0.2021 & 0.4181\\
1729 & H4 replay & 0.6492 & 0.1600 & 0.4046\\
1729 & H1 aligned & 0.6415 & 0.1714 & 0.4065\\
1729 & H2 aligned & 0.6188 & 0.1927 & 0.4057\\
1729 & DeepSets aligned & 0.6136 & 0.2044 & 0.4090\\
1729 & H4 aligned & 0.5902 & 0.1630 & 0.3766\\
2026 & H1 aligned & 0.6289 & 0.1962 & 0.4126\\
2026 & H2 aligned & 0.6482 & 0.1865 & 0.4173\\
2026 & DeepSets aligned & 0.5777 & 0.1266 & 0.3521\\
2026 & H4 aligned & 0.6084 & 0.1578 & 0.3831\\
\bottomrule\end{tabular}\end{table}

Table~\ref{tab:qa} reports answer F1. For primary seed 1729, aligned H4 changes mean F1 from 0.4181 to 0.3766. Its difference from MMR is $-0.0537$ and from matched static replay is $-0.0281$. For seed 2026, H4 changes from 0.4304 to 0.3831; aligned H1, H2 and DeepSets score 0.4126, 0.4173 and 0.3521. Every method retains all 128 exposed questions. These descriptive comparisons do not establish independent confirmation.


\paragraph{Exact selection does not imply acceleration.}
The original exact granular-ball and k-means indexes matched Flat on 1,476
natural comparisons across 738 questions. In the subsequent fixed
32+32-query timing comparison, Flat took 0.335/0.546 seconds per question,
granular balls 0.426/0.653, and k-means 0.423/0.644. Granular balls reduced
candidate inner products by about 66\%/44\%, but increased token-feasibility
work and retained index overhead. These measurements include per-query
construction and network work; some ran alongside GPU reading, so they are
not isolated hardware benchmarks. We preserve this no-speedup finding and
do not repeat index tuning in the present intervention.
The bounded continuation uses 1.25 GPU process-wall hours and 986 actual reader calls, including eight uncached checks. Identical generation identities are reused; the 1,536 logical rows are not independent samples or fresh calls.


\section{Discussion and Conclusion}
Evidence selection has multiple objectives that cannot be substituted for
one another. Annotation supervision makes set prediction inexpensive, but
does not certify answer utility. Search can expose a mismatch that sampled
set comparisons miss. Exact indexing can preserve a poor choice perfectly,
and fewer inner products can still cost more time. The bounded continuation
therefore evaluates actual selected support and a fixed reader rather than
declaring success from loss reduction alone. H2 improves slightly over its
own static checkpoint in both seeds; H1 improves in only the primary seed.
All aligned models remain below MMR in F1 point estimate, and H4 declines
in both seeds. These controls support a conditional analysis of
training alignment, not an exclusive advantage of higher interaction order.

These findings are bounded by exposed development data, document overlap,
incomplete support annotation, one reader, and 128-candidate retrieval.
They establish neither open-domain gains nor an exclusive granular-ball
benefit. Independent confirmation requires a prospective data boundary.

\paragraph{Reproducibility and disclosure.}
Code and aggregate results accompany the paper; source text and answers
are not redistributed. AI assistance covered implementation, analysis,
literature checking, and writing. Authorship, funding, interests, licensing,
and final submission require human confirmation.

\clearpage
\bibliographystyle{../shared/splncs04}
\bibliography{../shared/references,../shared/additions}
\end{document}
